% GENERATED FILE. Edit canonical lessons, then run npm run generate:books.
% canonical-source-hash: fnv1a64:afa1dce787adb6b1
% canonical-lessons: TA-C528-eluttalar, TA-C528-naval, TA-C528-kavitai, TA-C528-patippi, TA-C528-toku

\chapter{Books for Pleasure}
\label{ch:ta-books-for-pleasure}
\hlchaptermodality{528}

\begin{chapteropening}
\textbf{By the end of this chapter:} \emph{I can talk about writers, novels and poems.}

Complete the last lesson of chapter 528: I can talk about writers, novels and poems.
\end{chapteropening}

\section[\texorpdfstring{\ta{எழுத்தாளர்} (eḻuttāḷar)}{எழுத்தாளர் (eḻuttāḷar)}]{\ta{எழுத்தாளர்} (eḻuttāḷar) — a writer; an author}

\label{lesson:TA-C528-eluttalar}



\begin{quote}
Before the new one: say the Tamil for a club, then the Tamil for to support.
\end{quote}

\subsection*{You'll want to know: \ta{எழுத்தாளர்}}
\textbf{\ta{எழுத்தாளர்}} — \emph{eḻuttāḷar} — \textquotedblleft{}a writer; an author\textquotedblright{}.

\begin{quote}
\ta{அவர்} \ta{புகழ்பெற்ற} \ta{தமிழ்} \ta{எழுத்தாளர்}.
\end{quote}

\begin{quote}
\emph{avar pukaḻpeṟṟa tamiḻ eḻuttāḷar.}
\end{quote}

\begin{quote}
He is a well-known Tamil writer.
\end{quote}

Plural: \ta{எழுத்தாளர்கள்}, \emph{eḻuttāḷarkaḷ}.

\ta{எழுத்து} \textquotedblleft{}writing\textquotedblright{} + -\ta{ஆளர்} \textquotedblleft{}one who …\textquotedblright{}. Compare \ta{எழுத்தர்} (\emph{eḻuttar}), a clerk.

\begin{cousinweb}
\textbf{\ta{எழுத்தாளர்}} (\emph{eḻuttāḷar}) is native Dravidian: \ta{எழுத்து} (\emph{eḻuttu}), \textquotedblleft{}writing\textquotedblright{}, from \ta{எழுது}, \textquotedblleft{}to write\textquotedblright{} (DEDR 853), with -\ta{ஆளர்}, \textquotedblleft{}one who does it\textquotedblright{}.
\end{cousinweb}

\begin{grammarlens}[title={What you've built}]
One more word for describing what you lived and felt.
\end{grammarlens}

\subsection*{Your turn}
\begin{itemize}
\raggedright
  \item \emph{Say it:} \emph{eḻuttāḷar}
  \item \emph{Say it:} \emph{eḻuttāḷar}, once more
  \item \emph{Say it:} \emph{avar pukaḻpeṟṟa tamiḻ eḻuttāḷar.} — He is a well-known Tamil writer.
  \item \emph{Recall:} say the Tamil for a club, then the Tamil for to support, then say \emph{eḻuttāḷar} again
\end{itemize}

\begin{tcolorbox}[breakable,colback=teal!4,colframe=teal!35!black,title={Before you move on}]
What does \ta{எழுத்தாளர்} mean? (\textquotedblleft{}A writer; an author\textquotedblright{}.) Say it once more.
\end{tcolorbox}
\section[\texorpdfstring{\ta{நாவல்} (nāval)}{நாவல் (nāval)}]{\ta{நாவல்} (nāval) — a novel}

\label{lesson:TA-C528-naval}



\begin{quote}
Before the new one: say the Tamil for to support, then the Tamil for a writer.
\end{quote}

\subsection*{You'll want to know: \ta{நாவல்}}
\textbf{\ta{நாவல்}} — \emph{nāval} — \textquotedblleft{}a novel\textquotedblright{}.

\begin{quote}
\ta{இந்த} \ta{நாவலை} \ta{நான்} \ta{இரண்டு} \ta{நாட்களில்} \ta{படித்து} \ta{முடித்தேன்}.
\end{quote}

\begin{quote}
\emph{inta nāvalai nāṉ iraṇṭu nāṭkaḷil paṭittu muṭittēṉ.}
\end{quote}

\begin{quote}
I finished reading this novel in two days.
\end{quote}

Plural: \ta{நாவல்கள்}, \emph{nāvalkaḷ}.

Formal writing also has the Tamil word \ta{புதினம்} for a novel.

\begin{cousinweb}
\textbf{\ta{நாவல்}} (\emph{nāval}) is a loan from English \emph{novel}.
\end{cousinweb}

\begin{grammarlens}[title={What you've built}]
One more word for describing what you lived and felt.
\end{grammarlens}

\subsection*{Your turn}
\begin{itemize}
\raggedright
  \item \emph{Say it:} \emph{nāval}
  \item \emph{Say it:} \emph{nāval}, once more
  \item \emph{Say it:} \emph{inta nāvalai nāṉ iraṇṭu nāṭkaḷil paṭittu muṭittēṉ.} — I finished reading this novel in two days.
  \item \emph{Recall:} say the Tamil for to support, then the Tamil for a writer, then say \emph{nāval} again
\end{itemize}

\begin{tcolorbox}[breakable,colback=teal!4,colframe=teal!35!black,title={Before you move on}]
What does \ta{நாவல்} mean? (\textquotedblleft{}A novel\textquotedblright{}.) Say it once more.
\end{tcolorbox}
\section[\texorpdfstring{\ta{கவிதை} (kavitai)}{கவிதை (kavitai)}]{\ta{கவிதை} (kavitai) — a poem; poetry}

\label{lesson:TA-C528-kavitai}



\begin{quote}
Before the new one: say the Tamil for a writer, then the Tamil for a novel.
\end{quote}

\subsection*{You'll want to know: \ta{கவிதை}}
\textbf{\ta{கவிதை}} — \emph{kavitai} — \textquotedblleft{}a poem; poetry\textquotedblright{}.

\begin{quote}
\ta{அவள்} \ta{தன்} \ta{அம்மாவைப்} \ta{பற்றி} \ta{ஒரு} \ta{கவிதை} \ta{எழுதினாள்}.
\end{quote}

\begin{quote}
\emph{avaḷ taṉ ammāvaip paṟṟi oru kavitai eḻutiṉāḷ.}
\end{quote}

\begin{quote}
She wrote a poem about her mother.
\end{quote}

The plural is \ta{கவிதைகள்}. A poet is a \ta{கவிஞர்}.

\begin{cousinweb}
\textbf{\ta{கவிதை}} (\emph{kavitai}) is borrowed from Sanskrit \emph{kavitā}, \textquotedblleft{}poetry; a poem\textquotedblright{} (Monier-Williams), with the long final vowel written -\ta{ஐ}.
\end{cousinweb}

\begin{grammarlens}[title={What you've built}]
One more word for describing what you lived and felt.
\end{grammarlens}

\subsection*{Your turn}
\begin{itemize}
\raggedright
  \item \emph{Say it:} \emph{kavitai}
  \item \emph{Say it:} \emph{kavitai}, once more
  \item \emph{Say it:} \emph{avaḷ taṉ ammāvaip paṟṟi oru kavitai eḻutiṉāḷ.} — She wrote a poem about her mother.
  \item \emph{Recall:} say the Tamil for a writer, then the Tamil for a novel, then say \emph{kavitai} again
\end{itemize}

\begin{tcolorbox}[breakable,colback=teal!4,colframe=teal!35!black,title={Before you move on}]
What does \ta{கவிதை} mean? (\textquotedblleft{}A poem; poetry\textquotedblright{}.) Say it once more.
\end{tcolorbox}
\section[\texorpdfstring{\ta{பதிப்பி} (patippi)}{பதிப்பி (patippi)}]{\ta{பதிப்பி} (patippi) — to print and publish}

\label{lesson:TA-C528-patippi}



\begin{quote}
Before the new one: say the Tamil for a novel, then the Tamil for a poem.
\end{quote}

\subsection*{You'll want to know: \ta{பதிப்பி}}
\textbf{\ta{பதிப்பி}} — \emph{patippi} — \textquotedblleft{}to print and publish\textquotedblright{}.

\begin{quote}
\ta{ஒரு} \ta{சிறிய} \ta{நிறுவனம்} \ta{அவருடைய} \ta{முதல்} \ta{நூலைப்} \ta{பதிப்பித்தது}.
\end{quote}

\begin{quote}
\emph{oru ciṟiya niṟuvaṉam avaruṭaiya mutal nūlaip patippittatu.}
\end{quote}

\begin{quote}
A small firm published his first book.
\end{quote}

Past (I …): \ta{பதிப்பித்தேன்}, \emph{patippittēṉ}.

The book takes -\ta{ஐ}. The noun is \ta{பதிப்பு}, an edition.

\begin{cousinweb}
\textbf{\ta{பதிப்பி}} (\emph{patippi}) is native Dravidian: DEDR 3911 gives \ta{பதி} (\emph{pati}), \textquotedblleft{}to imprint, to stamp\textquotedblright{}, and \ta{பதிப்பு}, \textquotedblleft{}imprinting\textquotedblright{}; -\ta{ப்பி} makes the verb \textquotedblleft{}to print, to bring out\textquotedblright{}.
\end{cousinweb}

\begin{grammarlens}[title={What you've built}]
One more word for describing what you lived and felt.
\end{grammarlens}

\subsection*{Your turn}
\begin{itemize}
\raggedright
  \item \emph{Say it:} \emph{patippi}
  \item \emph{Say it:} \emph{patippi}, once more
  \item \emph{Say it:} \emph{oru ciṟiya niṟuvaṉam avaruṭaiya mutal nūlaip patippittatu.} — A small firm published his first book.
  \item \emph{Recall:} say the Tamil for a novel, then the Tamil for a poem, then say \emph{patippi} again
\end{itemize}

\begin{tcolorbox}[breakable,colback=teal!4,colframe=teal!35!black,title={Before you move on}]
What does \ta{பதிப்பி} mean? (\textquotedblleft{}To print and publish\textquotedblright{}.) Say it once more.
\end{tcolorbox}
\section[\texorpdfstring{\ta{தொகு} (toku)}{தொகு (toku)}]{\ta{தொகு} (toku) — to compile; to collect}

\label{lesson:TA-C528-toku}



\begin{quote}
Before the new one: say the Tamil for a poem, then the Tamil for to print and publish.
\end{quote}

\subsection*{You'll want to know: \ta{தொகு}}
\textbf{\ta{தொகு}} — \emph{toku} — \textquotedblleft{}to compile; to collect\textquotedblright{}.

\begin{quote}
\ta{அவர்} \ta{பழைய} \ta{நாட்டுப்புறப்} \ta{பாடல்களைத்} \ta{தொகுத்தார்}.
\end{quote}

\begin{quote}
\emph{avar paḻaiya nāṭṭuppuṟap pāṭalkaḷait tokuttār.}
\end{quote}

\begin{quote}
He compiled old folk songs.
\end{quote}

Past (I …): \ta{தொகுத்தேன்}, \emph{tokuttēṉ}.

What is compiled takes -\ta{ஐ}. The result is a \ta{தொகுப்பு}, a collection.

\begin{cousinweb}
\textbf{\ta{தொகு}} (\emph{toku}) is native Dravidian: DEDR 3476 gives \ta{தொகு} (\emph{toku}), \textquotedblleft{}to assemble, to collect, to gather together\textquotedblright{}.
\end{cousinweb}

\begin{grammarlens}[title={What you've built}]
One more word for describing what you lived and felt.
\end{grammarlens}

\subsection*{Your turn}
\begin{itemize}
\raggedright
  \item \emph{Say it:} \emph{toku}
  \item \emph{Say it:} \emph{toku}, once more
  \item \emph{Say it:} \emph{avar paḻaiya nāṭṭuppuṟap pāṭalkaḷait tokuttār.} — He compiled old folk songs.
  \item \emph{Recall:} say the Tamil for a poem, then the Tamil for to print and publish, then say \emph{toku} again
\end{itemize}

\begin{tcolorbox}[breakable,colback=teal!4,colframe=teal!35!black,title={Before you move on}]
What does \ta{தொகு} mean? (\textquotedblleft{}To compile; to collect\textquotedblright{}.) Say it once more.
\end{tcolorbox}
